% Lösungen Flächeninhalt und Volumen im Raum (zusammenbau v0.4, Heft)
\einheitenkopf{Flächeninhalt und Volumen im Raum}

\erg{89}{a) Höhe zu PQ – sie steht senkrecht auf der Grundseite \quad b) $|PQ| = 5$, $|QR| = 10$, $A = \tfrac{1}{2} \cdot 5 \cdot 10 = 25$ \quad c) $S_x(4 | 0 | 0)$, $S_y(0 | -6 | 0)$, Katheten 4 und 6: $A = \tfrac{1}{2} \cdot 4 \cdot 6 = 12$}
\erg{90}{a) $|OQ| = 6$, $M(1 | 2 | 2)$, $|MR| = 6$, $A = \tfrac{1}{2} \cdot 6 \cdot 6 = 18$ \quad b) $|OP| = 10$, die Höhe ist $z$ (R liegt über der Basismitte): $\tfrac{1}{2} \cdot 10 \cdot z = 40$, $z = 8$ \quad c) $T(-4 | 1 | 1)$; $RT$ ist parallel zu $PQ$ und dreimal so lang, gleiche Höhe: Trapez mit $\tfrac{1}{2} \cdot (a + 3a) \cdot h = 2ah$, Verhältnis $= \tfrac{1}{4}$}
\erg{91}{a) halbes Diagonalenprodukt – es ist eine Raute \quad b) $\overrightarrow{PQ} \cdot \overrightarrow{QR} = 0$, $|PQ| = 5$, $|QR| = 10$, $A = 5 \cdot 10 = 50$ \quad c) $|PR| = 6$, $|QS| = 12$, $A = \tfrac{1}{2} \cdot 6 \cdot 12 = 36$}
\erg{92}{a) $|PQ| = \sqrt{36 + 64} = 10$, $A = 10^2 = 100$ \quad b) $|PR| = 6\sqrt{2}$, $|QS| = 2\sqrt{6}$, eine Raute $A = \tfrac{1}{2} \cdot 6\sqrt{2} \cdot 2\sqrt{6} = 12\sqrt{3} \approx 20{,}78\,\mathrm{m}^2$, Dach $\approx 83{,}14\,\mathrm{m}^2$ \quad c) $|PQ| = 8$, $|SR| = 4$, Seitenmitten $M_1(4 | 1 | 1)$ und $M_2(4 | 4 | 5)$, $h = |M_1 M_2| = 5$, $A = \tfrac{1}{2} \cdot (8 + 4) \cdot 5 = 30$ \quad d) $z = 0$ für $t = 2$: $M(6 | 8 | 0)$, $|OM| = 10$; vier rechtwinklige Dreiecke mit den Katheten 10: $A = 4 \cdot \tfrac{1}{2} \cdot 100 = 200$ \quad e) $P'(0 | 1 + \lambda | 0)$, $S'(0 | 9 - \lambda | 0)$; Trapez mit den parallelen Seiten 8 und $8 - 2\lambda$, Höhe 8: $\tfrac{1}{2} \cdot (16 - 2\lambda) \cdot 8 = 48$, $\lambda = 2$: $P'(0 | 3 | 0)$, $S'(0 | 7 | 0)$ \quad f) $P'(0 | 2 + \lambda | 0)$, $S'(0 | 12 - \lambda | 0)$; Trapez mit den parallelen Seiten 10 und $10 - 2\lambda$, Höhe 10: $\tfrac{1}{2} \cdot (20 - 2\lambda) \cdot 10 = 60$, $\lambda = 4$: $P'(0 | 6 | 0)$, $S'(0 | 8 | 0)$ \quad g) $|PQ| = 6$, $A_{PQRS} = 36$; $|KL| = \sqrt{18}$, $A_{KLTU} = 18$; $36 = 2 \cdot 18$ \quad h) $|PQ| = 8$, $A_{PQRS} = 64$; $|KL| = \sqrt{32}$, $A_{KLTU} = 32$; $64 = 2 \cdot 32$}
\erg{93}{a) Im Querschnitt ist die geneigte Seite Ankathete des Neigungswinkels, ihr Schatten die Hypotenuse: Länge $= 1 : \mathrm{cos}\, 40^\circ \approx 1{,}31\,\mathrm{m}$; die waagerechte Seite bleibt 1{,}6\,m: $A \approx 2{,}09\,\mathrm{m}^2$}
\erg{94}{a) Die waagerechte Seite wird in wahrer Länge abgebildet. Im Querschnitt senkrecht zu ihr ist die geneigte Seite eine Kathete eines rechtwinkligen Dreiecks, dessen Hypotenuse ihr Schatten auf dem Boden ist; die Hypotenuse ist länger. Also ist das Schattenrechteck größer als das Vordach.}
\erg{95}{a) mit Faktor ein Drittel – der Körper hat eine Spitze (Pyramide) \quad b) $G = 16$, $h = 6$, $V = \tfrac{1}{3} \cdot 16 \cdot 6 = 32$ \quad c) $G = 40$, $\tfrac{1}{3} \cdot 40 \cdot h = 100$, $h = 7{,}5$}
\erg{96}{a) OS steht senkrecht auf der Grundfläche, ist also die Höhe; $G = \tfrac{1}{2} \cdot 5 \cdot 6 = 15$; $\tfrac{1}{3} \cdot 15 \cdot |OS| = 40$, $|OS| = 8$ \quad b) $G = 36$, $\tfrac{1}{3} \cdot 36 \cdot h = 96$, $h = 8$; $z = 2 + 8 = 10$ (oder $z = 2 - 8 = -6$) \quad c) Diagonalen $|PR| = 7$, $|QT| = 6$, $G = \tfrac{1}{2} \cdot 7 \cdot 6 = 21$, $V = \tfrac{1}{3} \cdot 21 \cdot 5 = 35$ \quad d) Diagonalen 6 und 14, $G = \tfrac{1}{2} \cdot 6 \cdot 14 = 42$, $h = 5$, $V = 42 \cdot 5 = 210$ \quad e) Umfang $= 30$, $30 \cdot h = 150$, $h = 5$ \quad f) $|QR| = \sqrt{169 - 25} = 12\,\mathrm{cm}$, $G = \tfrac{1}{2} \cdot 5 \cdot 12 = 30\,\mathrm{cm}^2$, $V = \tfrac{1}{3} \cdot 30 \cdot 9 = 90\,\mathrm{cm}^3$ \quad g) $|QR| = \sqrt{100 - 36} = 8\,\mathrm{cm}$, $G = \tfrac{1}{2} \cdot 6 \cdot 8 = 24\,\mathrm{cm}^2$, $V = \tfrac{1}{3} \cdot 24 \cdot 5 = 40\,\mathrm{cm}^3$}
\erg{97}{a) als Summe: Quader plus Dachprisma \quad b) Pyramide $PQUT$: $V = \tfrac{1}{3} \cdot 18 \cdot 8 = 48$; über der Deckfläche ragt die ähnliche Pyramide mit den Kanten 3, 3 und 4 heraus: $V = \tfrac{1}{3} \cdot 4{,}5 \cdot 4 = 6$; Teilkörper $= 48 - 6 = 42$ \quad c) In der Ebene gilt $z = \tfrac{y}{2}$; der kleinere Teilkörper ist ein Prisma der Länge 8 mit dreieckigem Querschnitt (Katheten 8 und 4): $V = \tfrac{1}{2} \cdot 8 \cdot 4 \cdot 8 = 128$}
\erg{98}{a) Pyramide $PQUT$: $V = \tfrac{1}{3} \cdot 18 \cdot 10 = 60$; über der Deckfläche ragt die ähnliche Pyramide mit den Kanten 3, 3 und 5 heraus: $V = \tfrac{1}{3} \cdot 4{,}5 \cdot 5 = 7{,}5$; Teilkörper $= 60 - 7{,}5 = 52{,}5$ \quad b) Die verlängerten Seitenkanten treffen sich in der Spitze der ganzen Pyramide; die Deckfläche halbiert die Seite, also liegt die Spitze 6\,m über der Deckfläche und 12\,m über der Grundfläche. Erster Summand: ganze Pyramide (Grundfläche $36\,\mathrm{m}^2$, Höhe 12\,m), zweiter: abgeschnittene Spitze (Grundfläche $9\,\mathrm{m}^2$, Höhe 6\,m); $V = 144 - 18 = 126\,\mathrm{m}^3$ \quad c) Erster Summand: Pyramide mit der Spitze auf der $z$-Achse bei 9 über dem rechtwinkligen Dreieck mit den Katheten 9 auf der $x$- und $y$-Achse; zweiter: die ähnliche Pyramide $OKLN$, deren Katheten und Höhe $k$ sind; $K$, $L$, $N$ liegen bei 4 auf den Achsen: $k = 4$ \quad d) $V_Q = G \cdot h$, $V_P = \tfrac{1}{3} \cdot G \cdot h$, $V_Q : V_P = 3$; Zuwachs $= 200\,\%$ \quad e) Kegel mit $r = 6$, $h = 8$, Mantellinie $s = 10$; $O = \pi \cdot 36 + \pi \cdot 6 \cdot 10 = 96\pi \approx 301{,}59$ \quad f) $10\,\mathrm{l} = 10\,000\,\mathrm{cm}^3$, $h = 10\,000 : (\pi \cdot 10^2) \approx 31{,}83\,\mathrm{cm}$ – höher als 30\,cm, sie passt nicht}
\erg{99}{a) $V(k) = \tfrac{1}{3} \cdot \tfrac{1}{2} \cdot (7 - k)(1 + k) \cdot 3 = \tfrac{1}{2} \cdot (7 - k)(1 + k)$: nach unten geöffnete Parabel mit den Nullstellen $-1$ und 7, das Maximum liegt in der Mitte: $k = 3$ \quad b) $V(k) = \tfrac{1}{3} \cdot \tfrac{1}{2} \cdot (4 + k)(8 - k) \cdot 9 = 1{,}5 \cdot (4 + k)(8 - k)$: nach unten geöffnete Parabel mit den Nullstellen $-4$ und 8, das Maximum liegt in der Mitte: $k = 2$ \quad c) $V_1 = \tfrac{1}{3} \cdot 36 \cdot (4 - k)$, $V_2 = \tfrac{1}{3} \cdot 36 \cdot (8 - (4 + k))$, zusammen $24 \cdot (4 - k) = 72$, $k = 1$ \quad d) $V_1 = \tfrac{1}{3} \cdot 20 \cdot (5 - k)$, $V_2 = \tfrac{1}{3} \cdot 20 \cdot (10 - (5 + k))$, zusammen $\tfrac{40}{3} \cdot (5 - k) = 40$, $k = 2$ \quad e) Graph II – $V(t) = 72 - \tfrac{1}{3} \cdot 12 \cdot (6 - 2t) = 48 + 8t$ ist linear in $t$, weil nur die Höhe der herausgenommenen Pyramide linear von $t$ abhängt; der Graph ist eine Gerade. \quad f) Graph I – $V(t) = 108 - \tfrac{1}{3} \cdot 12 \cdot (9 - 3t) = 72 + 12t$ ist linear in $t$, weil nur die Höhe der herausgenommenen Pyramide linear von $t$ abhängt; der Graph ist eine Gerade.}
